\documentclass[12pt,letterpaper]{article}
\usepackage[margin=1in]{geometry}
\usepackage{amsmath,amssymb,amsfonts}
\usepackage{enumitem}
\usepackage{booktabs}
\usepackage{hyperref}

% ---- macros consistent with main papers ----
\newcommand{\rhoa}{\rho_{\!A}}
\newcommand{\Ka}{K_{\!A}}
\newcommand{\rloop}{r_{\mathrm{loop}}}

\title{Investigation: Electron as a Vortex Ring\\
with a Bound Aether Core}
\author{Erik Otto}
\date{April 2026 --- Working Document}

\begin{document}
\maketitle

\begin{abstract}
We investigate a formation model for the electron in which a quantized
vortex ring, created in the high-energy early universe, accumulates a bound
Aether core through the Dynamic Casimir Effect (DCE).  The bound ring grows
until the core-boundary velocity drops below the DCE permanence threshold,
at which point only transient binding (gravity) occurs and the electron mass
stabilizes.  Using the theory's fundamental parameters ($m_A \approx
500$--$850$~MeV/$c^2$, $\rhoa \approx 4 \times 10^{11}$--$10^{12}$~kg/m$^3$),
the self-consistent vortex equation predicts $m_e \approx 0.4$--$0.6$~MeV,
bracketing the observed $0.511$~MeV.  Including corrections for the
core energy ($+5$--$14\%$, model-dependent) and the Bernoulli density
deficit near the vortex core ($-1.7\%$), the prediction matches $0.511$~MeV exactly for
$m_A \approx 628$--$689$~MeV/$c^2$, well within the independently
constrained convergence zone.
\end{abstract}


% ================================================================
\section{Formation Narrative}
\label{sec:narrative}
% ================================================================

The electron forms through the following sequence in the early universe:

\begin{enumerate}
  \item \textbf{Vortex ring creation.}  In the high-energy early universe,
    turbulent Aether flows create quantized vortex rings with circulation
    $\kappa_e = h/(2m_A)$---the minimum (spin-$\tfrac{1}{2}$) quantum.
    Initially these are pure fluid vortices: no bound material, no core
    structure beyond the healing-length depletion.

  \item \textbf{DCE creates matter.}  The circulation velocity at the
    healing length is $v(\xi) = c/\sqrt{2} \approx 0.707\,c$
    (Eq.~\ref{eq:v-healing}).  This exceeds the DCE permanence threshold,
    so the rapidly-circulating boundary continuously converts vacuum
    fluctuations into permanent bound Aether.

  \item \textbf{Matter migrates inward.}  The Bernoulli pressure in a vortex
    is lowest at the core (Eq.~\ref{eq:bernoulli}).  Bound material, no
    longer part of the superfluid flow, is pushed inward by the pressure
    gradient and accumulates along the vortex tube axis---forming a toroidal
    ring of bound Aether tracing the vortex's major radius $\rloop$.

  \item \textbf{Ring grows.}  As more bound material falls to the center,
    the ring expands radially outward from the tube axis.  The effective
    core boundary moves from $\xi$ to a larger radius $r_{\mathrm{eq}}$.

  \item \textbf{Boundary velocity drops.}  The circulation $\kappa_e$ is
    quantized and fixed, but the velocity at the boundary is
    $v(r) = \kappa_e/(2\pi r)$.  As $r_{\mathrm{eq}}$ increases, the
    boundary velocity decreases.

  \item \textbf{Permanence threshold reached.}  When $v(r_{\mathrm{eq}})$
    drops to $v_{\mathrm{perm}}$---the minimum velocity for permanent
    DCE binding---the vortex can no longer create permanent mass.  Only
    transient bind--unbind cycling occurs at the boundary, producing
    the pressure asymmetry that manifests as gravity (Paper~3).

  \item \textbf{Mass stabilizes.}  The electron's mass is the sum of the
    circulation kinetic energy and the bound ring's rest mass at this
    equilibrium point.
\end{enumerate}

This is a \emph{self-regulating} process: the creation of mass enlarges the
core, which reduces the boundary velocity, which eventually shuts off mass
creation.  The electron mass is not a free parameter but an equilibrium
outcome determined by medium properties.

\textbf{Energy source.}  Where does the ring's rest-mass energy come from?
The DCE mechanism converts energy from the vortex's circulation into bound
material---the moving boundary does work on the vacuum fluctuations, and
the energy is drawn from the kinetic energy of the flow.  As the ring
grows, the circulation KE decreases by the same amount that the ring mass
increases.  The total energy (KE + ring mass) is therefore approximately
\emph{conserved} during the formation process: $m_e \approx
E_{\mathrm{KE,initial}} = E_{\mathrm{KE,final}} + m_{\mathrm{ring}}\,c^2$.
The electron's final mass is set by the initial vortex energy at creation,
not by the partitioning between KE and ring mass.


% ================================================================
\section{Key Equations}
\label{sec:equations}
% ================================================================

% ----------------------------------------------------------------
\subsection{Circulation Velocity}
% ----------------------------------------------------------------

For a quantized vortex with circulation $\kappa_e = h/(2m_A) = \pi\hbar/m_A$,
the flow velocity at distance $r$ from the tube axis is:
%
\begin{equation}
  v(r) = \frac{\kappa_e}{2\pi\,r}
  \label{eq:v-of-r}
\end{equation}

At the healing length $\xi = \hbar/(m_A c\sqrt{2})$:
%
\begin{equation}
  v(\xi) = \frac{\pi\hbar/m_A}{2\pi\,\hbar/(m_A c\sqrt{2})}
         = \frac{c}{\sqrt{2}} \approx 0.707\,c
  \label{eq:v-healing}
\end{equation}

% ----------------------------------------------------------------
\subsection{Bernoulli Pressure}
% ----------------------------------------------------------------

The pressure distribution in the vortex flow:
%
\begin{equation}
  P(r) = P_\infty - \tfrac{1}{2}\,\rhoa\,v^2(r)
       = P_\infty - \frac{\rhoa\,\kappa_e^2}{8\pi^2\,r^2}
  \label{eq:bernoulli}
\end{equation}

Pressure is lowest at the core center ($r \to 0$) and increases outward.
Bound material (not participating in superfluid flow) experiences a net
inward force from the surrounding higher-pressure fluid.

% ----------------------------------------------------------------
\subsection{Equilibrium Condition}
% ----------------------------------------------------------------

The bound ring grows until the boundary velocity reaches the permanence
threshold:
%
\begin{equation}
  v(r_{\mathrm{eq}}) = v_{\mathrm{perm}}
  \quad\Longrightarrow\quad
  r_{\mathrm{eq}} = \frac{\kappa_e}{2\pi\,v_{\mathrm{perm}}}
  \label{eq:req}
\end{equation}

For $v_{\mathrm{perm}} < c/\sqrt{2}$, the ring extends beyond the healing
length: $r_{\mathrm{eq}} > \xi$.

% ----------------------------------------------------------------
\subsection{Core Energy: Already Accounted For}
\label{sec:core-energy}
% ----------------------------------------------------------------

An important subtlety: does the Donnelly formula already include the core's
energy?

The Donnelly formula computes only the \emph{kinetic energy of the
circulating flow outside the core}.  In a standard GPE vortex, the total
energy also includes the core's gradient and potential energy---the cost of
depleting the condensate.  This core energy is of order $1/\ln(\rloop/\xi)$
of the kinetic energy.  For the electron's aspect ratio, this is
$\sim\!5$--$14\%$ depending on the condensate potential (the exact value for
the logarithmic potential used in this theory has not been computed from first
principles).

In the bound-core model, the core region is filled with bound material
instead of depleted vacuum.  The bound ring's rest mass \emph{replaces} the
gradient and potential energy of the depleted core---it is the physical
realization of the same core energy, not an additional contribution on top of
it.  The total electron energy is therefore:
%
\begin{equation}
  m_e\,c^2 \;\approx\; E_{\mathrm{KE}}(\text{circulation})
  + E_{\mathrm{core}}
  \label{eq:me-general}
\end{equation}
%
where $E_{\mathrm{core}}$ exists in both models (depleted or bound) and is
$\sim\!5$--$14\%$ of $E_{\mathrm{KE}}$.  In the depleted model it is
gradient + potential energy; in the bound model it is the ring's rest mass.
The Donnelly formula (which captures only $E_{\mathrm{KE}}$) undercounts
the total by this amount regardless of core interpretation.

% ----------------------------------------------------------------
\subsection{Bernoulli Density Deficit Near the Core}
\label{sec:bernoulli-density}
% ----------------------------------------------------------------

The Donnelly formula assumes constant density $\rhoa$ everywhere.  But the
vortex circulation itself produces a Bernoulli density deficit (fixes.txt,
Fix~4.1: ``venturi-like deficit from driven spin energy''):
%
\begin{equation}
  \frac{\delta\rho(r)}{\rhoa} \approx -\frac{v^2(r)}{2\,c^2}
  = -\frac{\kappa_e^2}{8\pi^2\,c^2\,r^2}
  \label{eq:density-deficit}
\end{equation}

At the healing length ($r = \xi$, $v = c/\sqrt{2}$), this gives
$\delta\rho/\rhoa \approx -25\%$---not negligible.  The density drops
rapidly with distance:

\medskip
\begin{center}
\begin{tabular}{@{}ccc@{}}
\toprule
$r/\xi$ & $v/c$ & $\delta\rho/\rhoa$ \\
\midrule
1.0 & 0.707 & $-25\%$ \\
1.5 & 0.471 & $-11\%$ \\
2.0 & 0.354 & $-6\%$ \\
5.0 & 0.141 & $-1\%$ \\
\bottomrule
\end{tabular}
\end{center}
\medskip

The density deficit reduces the kinetic energy because
$E_{\mathrm{KE}} \propto \rho(r)\,v^2(r)$.  The correction is
$\delta E = \int\!\tfrac{1}{2}\,(-\delta\rho)\,v^2\,dV$, giving:
%
\begin{equation}
  \frac{\delta E_{\mathrm{KE}}}{E_{\mathrm{KE}}} \approx
  -\frac{1}{8\,\ln(\rloop/\xi)}
  \approx -1.7\%
  \label{eq:ke-correction}
\end{equation}
%
(The factor $1/8$ arises because $\kappa_e^2/(16\pi^2 c^2 \xi^2) = 1/8$
for $v(\xi) = c/\sqrt{2}$.)  Note: $\delta\rho/\rhoa$ is a genuine density
change (not just a pressure change with constant $\rhoa$) because the vortex
circulation is \emph{driven} flow, not free-fall---the persistent spin energy
maintains the flow against Bernoulli pressure, creating a venturi-like deficit.
This is physically distinct from gravitational inflow, where free-fall dynamics
ensure $\Delta\rho = 0$ exactly (Paper~2).  For the logarithmic equation of
state, the exact compressible Bernoulli gives $\rho(r) = \rhoa\,e^{-v^2/(2c^2)}$
rather than the linearized $\rhoa\,[1 - v^2/(2c^2)]$; the integrated KE
correction differs by less than $3\%$ between exact and linear forms, so the
linearized result $-1.7\%$ is adequate.

The Donnelly formula itself is non-relativistic.  In the Aether framework,
relativistic effects emerge from the medium at $O(v^2/c^2)$: time dilation
from wave path geometry, length contraction from density changes (Paper~2).
The Bernoulli density correction is the leading such correction to the vortex
energy; higher-order terms are $O(v^4/c^4) \sim 0.4\%$, negligible at this
level of precision.

% ----------------------------------------------------------------
\subsection{Electron Mass}
% ----------------------------------------------------------------

Combining all contributions:
%
\begin{equation}
  \boxed{m_e\,c^2 \;\approx\;
  \underbrace{\tfrac{1}{2}\,\rhoa\,\kappa_e^2\,\rloop
  \left[\ln\!\left(\frac{8\,\rloop}{r_{\mathrm{eq}}}\right)
  - \tfrac{1}{2}\right]}_{\displaystyle E_{\mathrm{KE}}\;\text{(circulation)}}
  \;\times\;\underbrace{(1 - 0.017)}_{\text{Bernoulli}}
  \;+\; \underbrace{E_{\mathrm{core}}}_{\displaystyle\sim\!5\text{--}14\%\text{ of }E_{\mathrm{KE}}}}
  \label{eq:me-total}
\end{equation}

The first term is the Donnelly formula corrected for the Bernoulli density
deficit ($-1.7\%$).  The second is the core energy ($5$--$14\%$ of
$E_{\mathrm{KE}}$), which in the bound-core model is the ring's rest mass.
The net correction from the bare Donnelly formula is:
%
\begin{equation}
  m_e\,c^2 \approx E_{\mathrm{KE}}^{\mathrm{Donnelly}} \times
  (1 - 0.017 + f_{\mathrm{core}})
\end{equation}
%
where $f_{\mathrm{core}} \approx 0.05$--$0.14$ is the core energy fraction
(model-dependent).  The net correction ranges from $+3\%$ to $+12\%$
above the bare Donnelly value.

The loop radius is set by the quantum minimum:
$\rloop = \hbar/(m_e\,c)$ (the reduced Compton wavelength).


% ================================================================
\section{Mass Estimate}
\label{sec:estimate}
% ================================================================

% ----------------------------------------------------------------
\subsection{Circulation KE as a Lower Bound}
% ----------------------------------------------------------------

As a first step, ignore the ring mass and set $m_e c^2 = E_{\mathrm{KE}}$
alone.  Using the self-consistent condition $\rloop = \hbar/(m_e c)$ and
setting $r_{\mathrm{eq}} = \xi$:
%
\begin{equation}
  m_e^2\,c^3 = \tfrac{1}{2}\,\rhoa\,\kappa_e^2\,\hbar
  \left[\ln\!\left(\frac{8\,\rloop}{r_{\mathrm{eq}}}\right) - \tfrac{1}{2}\right]
  \label{eq:self-consistent}
\end{equation}
%
\begin{equation}
  m_e^2 = \frac{\pi^2\,\rhoa\,\hbar^3}{2\,m_A^2\,c^3}
  \left[\ln\!\left(\frac{8\sqrt{2}\,m_A}{m_e}\right) - \tfrac{1}{2}\right]
  \label{eq:self-consistent-explicit}
\end{equation}

This gives a \emph{lower bound} on $m_e$ (since it ignores the ring's mass).
Solving iteratively for $\rhoa = 5 \times 10^{11}$~kg/m$^3$:

\medskip
\begin{center}
\begin{tabular}{@{}ccccc@{}}
\toprule
$m_A$ & $\xi$ & $\rloop$ & $m_e$ (KE only) & vs.\ observed \\
(MeV/$c^2$) & (fm) & (fm) & (MeV) & \\
\midrule
500 & 0.28 & 325 & 0.61 & $+19\%$ \\
550 & 0.25 & 354 & 0.56 & $+9\%$ \\
600 & 0.23 & 383 & 0.515 & $+1\%$ \\
650 & 0.21 & 411 & 0.48 & $-6\%$ \\
700 & 0.20 & 439 & 0.45 & $-12\%$ \\
800 & 0.17 & 495 & 0.40 & $-22\%$ \\
\bottomrule
\end{tabular}
\end{center}
\medskip

The bare KE-only equation gives $m_e = 0.515$~MeV at $m_A = 600$.
Applying the net correction from
Sections~\ref{sec:core-energy}--\ref{sec:bernoulli-density} ($+3\%$ to
$+12\%$, depending on core energy) increases the prediction.  Matching
$0.511$~MeV requires $m_A$ above $600$~MeV (see below).

% ----------------------------------------------------------------
\subsection{Full Prediction with Energy Conservation}
% ----------------------------------------------------------------

Since DCE converts circulation KE into ring mass (energy conserved),
the total electron mass is set by the initial vortex energy, not by the
ring's size.  The self-consistent equation remains
Eq.~(\ref{eq:self-consistent-explicit})---but now $m_e$ represents the
\emph{total} energy (KE + ring), and the Donnelly formula gives only the
KE portion.

Including the core energy correction and Bernoulli deficit from
Sections~\ref{sec:core-energy}--\ref{sec:bernoulli-density}, the corrected
prediction depends on the core energy fraction $f_{\mathrm{core}}$:

\medskip
\begin{center}
\begin{tabular}{@{}ccccc@{}}
\toprule
$f_{\mathrm{core}}$ & Net correction & Best-fit $m_A$ & KE only (MeV) & Corrected (MeV) \\
\midrule
$5\%$ & $+3\%$ & 628~MeV & 0.495 & 0.511 \\
$8\%$ & $+6\%$ & 649~MeV & 0.480 & 0.510 \\
$11\%$ & $+9\%$ & 669~MeV & 0.468 & 0.512 \\
$14\%$ & $+12\%$ & 689~MeV & 0.455 & 0.511 \\
\bottomrule
\end{tabular}
\end{center}
\medskip

For any core energy in the $5$--$14\%$ range, there exists an $m_A$ in the
convergence zone ($500$--$850$~MeV/$c^2$) that matches the observed mass:
%
\begin{equation}
  \boxed{m_e \approx 0.511\;\text{MeV} \quad\text{(predicted, }m_A \approx
  628\text{--}689\;\text{MeV)}
  \qquad\text{vs.}\qquad
  m_e = 0.511\;\text{MeV} \quad\text{(observed)}}
  \label{eq:result}
\end{equation}

The match is exact for $m_A$ values well within the convergence zone.  The
Bernoulli density deficit ($-1.7\%$) partially offsets the core energy
($+5$--$14\%$), and the remaining net correction ($+3$--$12\%$) is absorbed
by $m_A$, which shifts from $600$~MeV (bare Donnelly) to $628$--$689$~MeV
(corrected).

% ----------------------------------------------------------------
\subsection{Effect of Threshold Velocity}
% ----------------------------------------------------------------

Under energy conservation, the total electron mass is fixed by the initial
vortex energy at creation.  The threshold velocity $v_{\mathrm{perm}}$
determines only the \emph{partitioning} between circulation KE and ring
rest mass, not the total.  As $v_{\mathrm{perm}}$ decreases (ring grows),
more KE converts to ring mass, but $m_e$ stays constant.

The ring mass at equilibrium is:
%
\begin{equation}
  m_{\mathrm{ring}}\,c^2 = f_{\mathrm{core}} \times m_e\,c^2
  \label{eq:mring}
\end{equation}
%
where $f_{\mathrm{core}} \approx 5$--$14\%$ (the fraction of total energy
in the core, see Section~\ref{sec:core-energy}).  This is a small fraction:
the electron remains circulation-dominated at equilibrium.


% ================================================================
\section{Physical Structure at Equilibrium}
\label{sec:structure}
% ================================================================

% ----------------------------------------------------------------
\subsection{Compression Ratio}
\label{sec:compression}
% ----------------------------------------------------------------

Even if the entire electron mass were packed into the core ring (the extreme
$f = 1$ case), the compression ratio would be remarkably low:
%
\begin{equation}
  X = \frac{\rho_{\mathrm{core}}}{\rhoa}
    = \frac{m_e}{\rhoa\,\cdot\,2\pi^2\,\rloop\,\xi^2}
    \approx 4\text{--}5
  \label{eq:compression}
\end{equation}

(varying from $\sim\!4.4$ at $m_A = 600$ to $\sim\!4.9$ at $m_A = 630$).
Only $\sim\!5$ times ambient density even in the extreme case---far below
the $\sim\!7 \times 10^7$ compression in QCD bridges.  This reflects the
electron's small mass ($0.511$~MeV) spread over a large circumference
($2\pi\rloop \approx 2400$~fm).

For the core energy range $f_{\mathrm{core}} = 5$--$14\%$, the ring carries
only $0.03$--$0.07$~MeV.  Spread over the core volume
($\sim\!300$--$400$~fm$^3$), the ring density is well below
$\rhoa$---the core is sparsely filled, not compressed.

% ----------------------------------------------------------------
\subsection{What Stops the Growth}
% ----------------------------------------------------------------

The self-regulation mechanism works as follows:

\begin{enumerate}[leftmargin=2em]
  \item DCE creates bound material at the core boundary where $v > v_{\mathrm{perm}}$
  \item Material falls inward, expanding the effective core from $\xi$ to $r_{\mathrm{eq}}$
  \item At $r_{\mathrm{eq}}$, the velocity has dropped: $v(r_{\mathrm{eq}}) < v(\xi)$
  \item When $v(r_{\mathrm{eq}}) = v_{\mathrm{perm}}$, permanent binding ceases
  \item Transient binding continues (producing gravity), but no new mass accumulates
\end{enumerate}

The quantized circulation $\kappa_e = h/(2m_A)$ is topologically fixed---it
cannot change.  Only the boundary radius changes as the ring grows, reducing
the velocity through the $v \propto 1/r$ dependence.


% ================================================================
\section{The DCE Permanence Threshold}
\label{sec:threshold}
% ================================================================

The permanence threshold $v_{\mathrm{perm}}$ is the single most important
unknown in this model.  It determines:
\begin{itemize}[leftmargin=2em]
  \item The equilibrium core size $r_{\mathrm{eq}} = \kappa_e/(2\pi v_{\mathrm{perm}})$
  \item The partitioning between circulation KE and bound ring mass
  \item The detailed structure of the electron
\end{itemize}

For $v_{\mathrm{perm}}$ near $c/\sqrt{2}$ (modest ring growth), the total
mass varies only weakly with $v_{\mathrm{perm}}$: the logarithmic KE
decrease partially cancels the ring mass increase.  For
$v_{\mathrm{perm}} \ll c$, however, the ring mass (scaling as
$r_{\mathrm{eq}}^2$) dominates the logarithmic KE decrease, and the total
mass grows without bound.  This constrains $v_{\mathrm{perm}}$ to be close
to $c/\sqrt{2}$.

The threshold must satisfy $v_{\mathrm{perm}} < c/\sqrt{2}$ (otherwise no
ring growth occurs) and $v_{\mathrm{perm}} > 0$ (otherwise the ring grows
without limit).  A natural estimate from the Schwinger pair-production
analogy---where the creation rate transitions from exponentially suppressed
to significant when the exponent $(c/v)^2 \sim \pi$---gives:
%
\begin{equation}
  v_{\mathrm{perm}} \sim \frac{c}{\sqrt{\pi}} \approx 0.56\,c
  \label{eq:threshold-estimate}
\end{equation}

This corresponds to $r_{\mathrm{eq}} \approx 1.26\,\xi$---a modest expansion
of the core beyond the healing length.

Quantifying $v_{\mathrm{perm}}$ from first principles (the DCE pair creation
rate in vortex geometry) is the key open problem identified in Paper~7.


% ================================================================
\section{Significance}
\label{sec:significance}
% ================================================================

\subsection{Electron Mass from Medium Properties}

The central result is Eq.~(\ref{eq:self-consistent-explicit}) combined with
Eq.~(\ref{eq:mring}): the electron mass is determined by $m_A$, $\rhoa$,
$c$, and the ring density $\alpha$---all properties of the Aether medium and
the equilibrium vortex structure, not free parameters of the electron.  The
mass arises from:

\begin{itemize}[leftmargin=2em]
  \item \textbf{Circulation quantization}: $\kappa_e = h/(2m_A)$ fixes the
    flow speed at each radius
  \item \textbf{Quantum minimum size}: $\rloop = \hbar/(m_e c)$ prevents the
    vortex ring from shrinking to zero
  \item \textbf{Self-consistency}: $m_e$ must equal the energy of a vortex
    ring of size $\rloop$, which itself depends on $m_e$
\end{itemize}

For $m_A \approx 628$--$689$~MeV (within the convergence zone), the
corrected prediction matches the observed $0.511$~MeV exactly.  The
inputs---$m_A$ and $\rhoa$---are independently constrained by QCD bridge
physics (Paper~5) and the healing-length convergence (Paper~7).  The core
energy fraction $f_{\mathrm{core}}$ is the remaining model-dependent
quantity; computing the vortex core structure in the logarithmic potential
would pin down both $f_{\mathrm{core}}$ and the best-fit $m_A$
simultaneously.

\subsection{Self-Regulating Mass}

The bound-core model explains \emph{why} the electron has the mass it does:
the DCE mechanism creates mass until it shuts itself off.  Every electron
reaches the same equilibrium because every minimum-circulation vortex ring
has the same $\kappa_e$, the same velocity profile, and therefore the same
threshold crossing point.  Mass universality follows from circulation
quantization.

\subsection{Open Tube / Closed Tube Taxonomy}

The bound electron core ring is a \emph{closed tube} of gently compressed
Aether, stabilized by the encircling vortex circulation.  The QCD bridge is
an \emph{open tube} of strongly compressed Aether, stabilized by quark
vortex endpoints.  Both are topologically-stabilized bound Aether structures
in the same medium, differing in topology (closed vs.\ open), compression
($\ll 1$ to $\sim\!5$ vs.\ $\sim\!7 \times 10^7$), and energy scale ($0.5$~MeV
vs.\ $\sim\!1$~GeV/fm).


% ================================================================
\section{Open Questions}
\label{sec:open}
% ================================================================

\begin{enumerate}[leftmargin=2em]
  \item \textbf{DCE permanence threshold.}  What is $v_{\mathrm{perm}}$?
    A first-principles DCE calculation in vortex geometry would determine the
    equilibrium core size and the KE/ring mass partitioning.

  \item \textbf{Muon and tau.}  Could higher-energy vortex configurations
    (multiply-wound cores, knotted rings) stabilize at different masses?
    The self-regulation mechanism would apply to any circulation quantum,
    potentially producing the lepton mass hierarchy.

  \item \textbf{Form factor.}  Electron scattering experiments probe to
    $\sim\!10^{-18}$~m with no structure seen, yet the vortex model predicts
    structure at $\rloop \sim 400$~fm.  How does the vortex ring's
    scattering cross-section mimic a point particle?

  \item \textbf{Charge origin.}  The bound core ring surface is a permanent
    DCE boundary.  Could the intense transient cycling at this surface be
    the microscopic mechanism that implements the radial Aether flow defining
    electric charge (an acknowledged gap in Paper~4)?

  \item \textbf{Cosmological formation.}  What conditions in the early
    universe create vortex rings with the minimum circulation quantum?  What
    determines the initial $\rloop$, and how does it relax to
    $\hbar/(m_e c)$?
\end{enumerate}


\end{document}
